\subsubsection{Frontend}

\begin{figure}[H]
	\centering 
	\includegraphics[width=1\textwidth]{SpecificaTecnica/Storico/StoricoFrontend}
	\caption{Diagramma delle classi per Storico Ordini: Frontend}
\end{figure}

\subsubsubsection{StoricoView}
\begin{figure}[H]
	\centering 
	\includegraphics[width=0.4\textwidth]{SpecificaTecnica/Storico/StoricoView}
	\caption{Storico Ordini: StoricoView}
\end{figure}

\paragraph{Descrizione}
Componente React che rappresenta la pagina dello storico degli ordini. \\
Recupera i dati tramite \texttt{useStoricoViewModel} e gestisce la visualizzazione della lista degli ordini con supporto alla paginazione e al filtraggio per data. \\
Se l'utente è amministratore, vengono mostrati gli ordini di tutti i clienti. \\
Utilizza i componenti \texttt{ComponentOrdineRow}, \texttt{ComponentDettaglioModal} e \texttt{ComponentPaginazione}.

\paragraph{Attributi}
\begin{itemize}
	\item \texttt{profileOpen: boolean} - indica se il pannello profilo utente è aperto.
	\item \texttt{filtroAperto: boolean} - indica se il pannello di filtraggio per data è aperto.
	\item \texttt{dataInizio: string} - data di inizio del filtro nel formato \texttt{YYYY-MM-DD}.
	\item \texttt{dataFine: string} - data di fine del filtro nel formato \texttt{YYYY-MM-DD}.
\end{itemize}

\paragraph{Metodi e funzioni}
\begin{itemize}
	\item \texttt{handleLogout(): void} - esegue il logout dell'utente chiamando l'endpoint \texttt{/api/auth/logout} e reindirizza alla pagina principale.
	\item \texttt{handleProfile(): void} - apre il pannello del profilo utente impostando \texttt{profileOpen} a \texttt{true}.
	\item \texttt{resetFiltro(): void} - azzera i campi \texttt{dataInizio} e \texttt{dataFine} e ricarica la prima pagina senza filtri.
	\item \texttt{ordiniFiltrati: Ordine[]} - restituisce la lista degli ordini filtrata in base alle date \texttt{dataInizio} e \texttt{dataFine} correntemente impostate.
\end{itemize}

\subsubsubsection{ComponentPaginazione}
\begin{figure}[H]
	\centering 
	\includegraphics[width=0.4\textwidth]{SpecificaTecnica/Storico/ComponentPaginazione}
	\caption{Storico Ordini: ComponentPaginazione}
\end{figure}

\paragraph{Descrizione}
Componente React che renderizza i controlli di paginazione per la lista degli ordini.\\
Riceve la pagina corrente, il numero totale di pagine e una callback per il cambio pagina.\\
Non viene renderizzato se il numero totale di pagine è inferiore o uguale a uno.

\paragraph{Attributi}
\begin{itemize}
	\item \texttt{pagina: number} - indice della pagina correntemente visualizzata.
	\item \texttt{totalePagine: number} - numero totale di pagine disponibili.
	\item \texttt{onCambia: (n: number) => void} - callback invocata al cambio di pagina, riceve il numero della pagina richiesta.
\end{itemize}

\paragraph{Metodi e funzioni}
\begin{itemize}
	\item Il componente non espone metodi propri.\\
	La navigazione tra le pagine è gestita interamente tramite la callback \texttt{onCambia}, invocata con \texttt{pagina~-~1} o \texttt{pagina~+~1} alla pressione dei pulsanti \texttt{Precedente} e \texttt{Successivo} rispettivamente.\\
	I pulsanti vengono disabilitati automaticamente quando ci si trova alla prima o all'ultima pagina.
\end{itemize}

\subsubsubsection{ComponentDettaglioModal}
\begin{figure}[H]
	\centering 
	\includegraphics[width=0.4\textwidth]{SpecificaTecnica/Storico/ComponentDettaglioModal}
	\caption{Storico Ordini: ComponentDettaglioModal}
\end{figure}

\paragraph{Descrizione}
Componente React che renderizza una finestra modale con il dettaglio di un ordine selezionato.\\
Mostra il codice ordine, la data, la lista dei prodotti con le rispettive quantità e, se l'utente è amministratore, il nome del cliente.\\
Per gli utenti non amministratori è presente un pulsante per duplicare l'ordine, con relativa richiesta di conferma.\\
La modale si chiude cliccando sul backdrop o sul pulsante di chiusura.

\paragraph{Attributi}
\begin{itemize}
	\item \texttt{ordine: Ordine} - oggetto contenente i dati dell'ordine da visualizzare.
	\item \texttt{isAdmin: boolean} - indica se l'utente corrente è un amministratore.
	\item \texttt{onChiudi: () => void} - callback invocata alla chiusura della modale.
	\item \texttt{onDuplica: (codice: string) => void} - callback invocata alla conferma della duplicazione dell'ordine.
	\item \texttt{erroreDuplica: string | null} - messaggio di errore da mostrare in caso di duplicazione fallita.
	\item \texttt{onRefresh: () => void} - callback invocata dopo una duplicazione avvenuta con successo per aggiornare la lista.
	\item \texttt{confermaDuplica: boolean} - indica se è in corso la richiesta di conferma della duplicazione.
	\item \texttt{duplicatoOk: boolean} - indica se la duplicazione è avvenuta con successo.
\end{itemize}

\paragraph{Metodi e funzioni}
\begin{itemize}
	\item \texttt{handleBackdropClick(): void} - chiude la modale quando l'utente clicca sul backdrop.
	\item \texttt{contenutoDuplica(): ReactNode} - restituisce il contenuto della sezione duplicazione in base allo stato corrente: il pulsante iniziale, la richiesta di conferma oppure il messaggio di successo.
\end{itemize}

\subsubsubsection{ComponentOrdineRow}
\begin{figure}[H]
	\centering 
	\includegraphics[width=0.4\textwidth]{SpecificaTecnica/Storico/ComponentOrdineRow}
	\caption{Storico Ordini: ComponentOrdineRow}
\end{figure}

\paragraph{Descrizione}
Componente React che renderizza una singola riga della tabella degli ordini.\\
Mostra il codice ordine, la data, il numero di prodotti e un pulsante per aprire il dettaglio.\\
Se l'utente è amministratore, viene mostrata anche la colonna con il nome del cliente.

\paragraph{Attributi}
\begin{itemize}
	\item \texttt{ordine: Ordine} - oggetto contenente i dati dell'ordine da visualizzare.
	\item \texttt{isAdmin: boolean} - indica se l'utente corrente è un amministratore.
	\item \texttt{onApriDettaglio: (o: Ordine) => void} - callback invocata alla pressione del pulsante di dettaglio, riceve l'ordine selezionato.
\end{itemize}

\paragraph{Metodi e funzioni}
\begin{itemize}
	\item Il componente non espone metodi propri.\\
	L'unica interazione è il pulsante \texttt{Dettaglio}, che invoca la callback \texttt{onApriDettaglio} passando l'ordine corrente.
\end{itemize}

\subsubsubsection{StoricoViewModel}
\begin{figure}[H]
	\centering 
	\includegraphics[width=0.4\textwidth]{SpecificaTecnica/Storico/StoricoViewModel}
	\caption{Storico Ordini: StoricoViewModel}
\end{figure}

\paragraph{Descrizione}
Hook React che funge da ViewModel per la gestione dello storico ordini.\\
Espone stato e azioni per caricare, visualizzare e duplicare ordini, adattando il comportamento al ruolo dell'utente autenticato.\\
Interagisce con \texttt{StoricoModel} per il recupero e la duplicazione dei dati tramite le funzioni \texttt{getStoricoCliente}, \texttt{getStoricoAdmin} e \texttt{duplicaOrdine}.

\paragraph{Attributi}
\begin{itemize}
	\item \texttt{ordini: Ordine[]} - lista degli ordini della pagina corrente.
	\item \texttt{pagina: number} - numero della pagina attualmente visualizzata.
	\item \texttt{totalePagine: number} - numero totale di pagine disponibili.
	\item \texttt{ordineScelto: Ordine | null} - ordine selezionato per il dettaglio, \texttt{null} se nessuno è selezionato.
	\item \texttt{loading: boolean} - indica se è in corso il caricamento degli ordini.
	\item \texttt{errore: string | null} - messaggio di errore del caricamento, \texttt{null} se assente.
	\item \texttt{erroreDuplica: string | null} - messaggio di errore della duplicazione, \texttt{null} se assente.
	\item \texttt{dataInizio: string} - data di inizio del filtro correntemente applicato.
	\item \texttt{dataFine: string} - data di fine del filtro correntemente applicato.
	\item \texttt{isAdmin: boolean} - indica se l'utente autenticato ha ruolo amministratore.
\end{itemize}

\paragraph{Metodi e funzioni}
\begin{itemize}
	\item \texttt{caricaPagina(n: number, di?: string, df?: string): Promise<void>} - carica la pagina \texttt{n} degli ordini applicando opzionalmente i filtri per data di inizio e di fine.\\
	Invoca \texttt{getStoricoAdmin} o \texttt{getStoricoCliente} in base al ruolo dell'utente.
	\item \texttt{apriDettaglio(ordine: Ordine): void} - imposta \texttt{ordineScelto} con l'ordine ricevuto, aprendo il dettaglio.
	\item \texttt{chiudiDettaglio(): void} - reimposta \texttt{ordineScelto} a \texttt{null}, chiudendo il dettaglio.
	\item \texttt{duplicaOrdine(codice: string): Promise<void>} - richiede la duplicazione dell'ordine identificato da \texttt{codice} tramite \texttt{apiDuplicaOrdine}.\\
	In caso di errore aggiorna \texttt{erroreDuplica} con il relativo messaggio.
\end{itemize}

\subsubsubsection{StoricoModel}
\begin{figure}[H]
	\centering 
	\includegraphics[width=0.4\textwidth]{SpecificaTecnica/Storico/StoricoModel}
	\caption{Storico Ordini: StoricoModel}
\end{figure}

\paragraph{Descrizione}
File TypeScript che definisce le interfacce dati e le funzioni di accesso alle API per lo storico ordini.\\
Espone i tipi \texttt{Ordine}, \texttt{Prodotto} e \texttt{StoricoPage} utilizzati internamente dai componenti, e i tipi \texttt{OrdineSchema}, \texttt{ProdottoSchema} e \texttt{StoricoPageSchema} che rispecchiano il formato restituito dal backend.\\
Le funzioni di accesso alle API gestiscono i parametri di paginazione e filtraggio per data.

\paragraph{Attributi}
\begin{itemize}
	\item \texttt{Prodotto} - interfaccia che rappresenta un prodotto all'interno di un ordine, con campi \texttt{nome}, \texttt{descrizione} e \texttt{quantita}.
	\item \texttt{Ordine} - interfaccia che rappresenta un ordine, con campi \texttt{codice\_ordine}, \texttt{numero\_ordine}, \texttt{data}, \texttt{username} opzionale e lista di \texttt{Prodotto}.
	\item \texttt{StoricoPage} - interfaccia che rappresenta una pagina di risultati, con campi \texttt{ordini}, \texttt{pagina\_corrente} e \texttt{totale\_pagine}.
	\item \texttt{ProdottoSchema} - interfaccia che rispecchia il formato prodotto restituito dal backend, con campi \texttt{nome} e \texttt{quantita}.
	\item \texttt{OrdineSchema} - interfaccia che rispecchia il formato ordine restituito dal backend, con campi \texttt{codice\_ordine}, \texttt{data}, \texttt{prodotti} e \texttt{cliente} opzionale.
	\item \texttt{StoricoPageSchema} - interfaccia che rispecchia il formato pagina restituito dal backend, con campi \texttt{ordini}, \texttt{pagina\_corrente} e \texttt{totale\_pagine}.
\end{itemize}

\paragraph{Metodi e funzioni}
\begin{itemize}
	\item \texttt{getStoricoCliente(pagina: number, perPagina: number, dataInizio?: string, dataFine?: string): Promise<StoricoPageSchema>} - recupera lo storico ordini dell'utente autenticato tramite \texttt{GET /api/storico/miei}, applicando opzionalmente i filtri per data.
	\item \texttt{getStoricoAdmin(pagina: number, perPagina: number, dataInizio?: string, dataFine?: string): Promise<StoricoPageSchema>} - recupera lo storico ordini di tutti gli utenti tramite \texttt{GET /api/storico/tutti}, applicando opzionalmente i filtri per data.
	\item \texttt{duplicaOrdine(codiceOrdine: string): Promise<void>} - duplica un ordine esistente tramite \texttt{POST /api/storico/duplica/\{codiceOrdine\}}.\\
	In caso di risposta non positiva lancia un errore con il messaggio restituito dal backend.
\end{itemize}